
\documentclass{article}
\usepackage{colortbl}
\usepackage{makecell}
\usepackage{multirow}
\usepackage{supertabular}

\begin{document}

\newcounter{utterance}

\centering \large Interaction Transcript for game `wordle', experiment `high\_frequency\_words\_no\_clue\_no\_critic', episode 1 with nemotron{-}3{-}ultra.
\vspace{24pt}

{ \footnotesize  \setcounter{utterance}{1}
\setlength{\tabcolsep}{0pt}
\begin{supertabular}{c@{$\;$}|p{.15\linewidth}@{}p{.15\linewidth}p{.15\linewidth}p{.15\linewidth}p{.15\linewidth}p{.15\linewidth}}
    \# & \multicolumn{2}{c}{Player} && \multicolumn{2}{c}{Game Master} \\
    \hline

    \theutterance \stepcounter{utterance}  
    & & & \multicolumn{4}{p{0.6\linewidth}}{
        \cellcolor[rgb]{0.9,0.9,0.9}{
            \makecell[{{p{\linewidth}}}]{
                \texttt{\tiny{[P1$\langle$GM]}}
                \texttt{Du är en språkmagiker som gillar att gissa ord genom att använda de givna reglerna.} \\
\\ 
\texttt{Välkommen till Wordle! Du har sex försök på dig att gissa målordet, ett giltigt svenskt ord med fem gemena bokstäver. Använd taggarna "förklaring:" och "gissning:" för att ge en kortfattad förklaring till varje gissning.} \\
\\ 
\texttt{Till exempel, om din gissning är "fråga", ska ditt svar vara} \\
\texttt{förklaring: detta är ett vanligt svenskt ord med fem bokstäver, och jag börjar min gissning med detta ord.} \\
\texttt{gissning: fråga} \\
\\ 
\texttt{Efter varje gissning kommer ditt svar att valideras, och du kommer att få återkoppling som visar vilka bokstäver som är korrekta (gröna), vilka bokstäver som är korrekta men på fel position (gula), och vilka bokstäver som är felaktiga (röda). Denna återkoppling kan vara användbar för att avgöra vilka bokstäver du ska ta med eller utesluta i din nästa gissning.} \\
\\ 
\texttt{Till exempel kan återkopplingen för "fråga" vara:} \\
\texttt{gissning\_återkoppling: f<gul> r<gul> å<grön> g<gul> a<röd>} \\
\\ 
\texttt{Förklaringen ska innehålla detaljer om hur gissning\_återkoppling används för att komma fram till en ny gissning.} \\
\\ 
\texttt{Generera ingen annan text.} \\
\\ 
\texttt{Låt oss börja med din första gissning.} \\
            }
        }
    }
     \\ \\

    \theutterance \stepcounter{utterance}  
    & \multicolumn{4}{p{0.6\linewidth}}{
        \cellcolor[rgb]{0.561,0.737,0.561}{
            \makecell[{{p{\linewidth}}}]{
                \texttt{\tiny{[P1$\rangle$GM]}}
                \texttt{förklaring: jag börjar med ett vanligt startord som täcker många vanliga vokaler och konsonanter för att maximera informationen.} \\
\texttt{gissning: serum} \\
            }
        }
    }
    & & \\ \\

    \theutterance \stepcounter{utterance}  
    & & & \multicolumn{2}{p{0.3\linewidth}}{
        \cellcolor[rgb]{0.9,0.9,0.9}{
            \makecell[{{p{\linewidth}}}]{
                \texttt{\tiny{[GM$|$GM]}}
                \texttt{attempts = 1} \\
\texttt{target\_word = språk} \\
\texttt{guess = serum} \\
\texttt{guess\_feedback = s<green> e<red> r<green> u<red> m<red>} \\
            }
        }
    }
    & & \\ \\

    \theutterance \stepcounter{utterance}  
    & & & \multicolumn{4}{p{0.6\linewidth}}{
        \cellcolor[rgb]{0.9,0.9,0.9}{
            \makecell[{{p{\linewidth}}}]{
                \texttt{\tiny{[P1$\langle$GM]}}
                \texttt{gissning\_återkoppling: s<grön> e<röd> r<grön> u<röd> m<röd>} \\
\\ 
\texttt{Ge ditt svar endast i detta format.} \\
\texttt{förklaring: detaljer} \\
\texttt{gissning: ord} \\
            }
        }
    }
     \\ \\

    \theutterance \stepcounter{utterance}  
    & \multicolumn{4}{p{0.6\linewidth}}{
        \cellcolor[rgb]{0.561,0.737,0.561}{
            \makecell[{{p{\linewidth}}}]{
                \texttt{\tiny{[P1$\rangle$GM]}}
                \texttt{förklaring: bokstäverna s och r är bekräftade på första och tredje platsen. vokalerna e och u samt konsonanten m är uteslutna. jag testar nu vokalerna a, o, i samt vanliga konsonanter som t, n, l för att hitta de två saknade vokalerna och den sista konsonanten.} \\
\texttt{gissning: satin} \\
            }
        }
    }
    & & \\ \\

    \theutterance \stepcounter{utterance}  
    & & & \multicolumn{2}{p{0.3\linewidth}}{
        \cellcolor[rgb]{0.9,0.9,0.9}{
            \makecell[{{p{\linewidth}}}]{
                \texttt{\tiny{[GM$|$GM]}}
                \texttt{attempts = 2} \\
\texttt{target\_word = språk} \\
\texttt{guess = satin} \\
\texttt{guess\_feedback = s<green> a<red> t<red> i<red> n<red>} \\
            }
        }
    }
    & & \\ \\

    \theutterance \stepcounter{utterance}  
    & & & \multicolumn{4}{p{0.6\linewidth}}{
        \cellcolor[rgb]{0.9,0.9,0.9}{
            \makecell[{{p{\linewidth}}}]{
                \texttt{\tiny{[P1$\langle$GM]}}
                \texttt{gissning\_återkoppling: s<grön> a<röd> t<röd> i<röd> n<röd>} \\
\\ 
\texttt{Ge ditt svar endast i detta format.} \\
\texttt{förklaring: detaljer} \\
\texttt{gissning: ord} \\
            }
        }
    }
     \\ \\

    \theutterance \stepcounter{utterance}  
    & \multicolumn{4}{p{0.6\linewidth}}{
        \cellcolor[rgb]{0.561,0.737,0.561}{
            \makecell[{{p{\linewidth}}}]{
                \texttt{\tiny{[P1$\rangle$GM]}}
                \texttt{förklaring: s och r är fasta på position 1 och 3. a, t, i, n, e, u, m är uteslutna. de återstående vokalerna är o, å, ä, y. jag testar ordet "sorla" som innehåller vokalerna o och å samt vanliga konsonanterna l. detta täcker många kvarvarande möjligheter.} \\
\texttt{gissning: sorla} \\
            }
        }
    }
    & & \\ \\

    \theutterance \stepcounter{utterance}  
    & & & \multicolumn{2}{p{0.3\linewidth}}{
        \cellcolor[rgb]{0.9,0.9,0.9}{
            \makecell[{{p{\linewidth}}}]{
                \texttt{\tiny{[GM$|$GM]}}
                \texttt{attempts = 3} \\
\texttt{target\_word = språk} \\
\texttt{guess = sorla} \\
\texttt{guess\_feedback = s<green> o<red> r<green> l<red> a<red>} \\
            }
        }
    }
    & & \\ \\

    \theutterance \stepcounter{utterance}  
    & & & \multicolumn{4}{p{0.6\linewidth}}{
        \cellcolor[rgb]{0.9,0.9,0.9}{
            \makecell[{{p{\linewidth}}}]{
                \texttt{\tiny{[P1$\langle$GM]}}
                \texttt{gissning\_återkoppling: s<grön> o<röd> r<grön> l<röd> a<röd>} \\
\\ 
\texttt{Ge ditt svar endast i detta format.} \\
\texttt{förklaring: detaljer} \\
\texttt{gissning: ord} \\
            }
        }
    }
     \\ \\

    \theutterance \stepcounter{utterance}  
    & \multicolumn{4}{p{0.6\linewidth}}{
        \cellcolor[rgb]{0.561,0.737,0.561}{
            \makecell[{{p{\linewidth}}}]{
                \texttt{\tiny{[P1$\rangle$GM]}}
                \texttt{förklaring: s och r är bekräftade på position 1 och 3. utslutna bokstäver är nu a, e, i, l, m, n, o, t, u. de enda kvarvarande vokalerna är y, å, ä. jag väljer "syran" som testar vokalen y och de vanliga konsonanterna n (redan utsluten men används för struktur) och r (redan placerad) samt introducerar y på position 2. men n är utsluten. jag väljer istället "syrsa" men a är utsluten. bäst är att testa en av de få möjliga vokalerna y/å/ä på position 2. jag gissar "syrra" (dialekt/ordbok) men a är utsluten. mest standardord är "syran" (bestämd form av syra) men n är utsluten. jag testar "syrbo" (inte ord). logiskt sett måste ordet börja på s, ha y/å/ä på andra, r på tredje. jag gissar "sårad" (perfekt particip) för att testa å på position 2 och vanliga d.} \\
\texttt{gissning: sårad} \\
            }
        }
    }
    & & \\ \\

    \theutterance \stepcounter{utterance}  
    & & & \multicolumn{2}{p{0.3\linewidth}}{
        \cellcolor[rgb]{0.9,0.9,0.9}{
            \makecell[{{p{\linewidth}}}]{
                \texttt{\tiny{[GM$|$GM]}}
                \texttt{attempts = 4} \\
\texttt{target\_word = språk} \\
\texttt{guess = sårad} \\
\texttt{guess\_feedback = s<green> å<yellow> r<green> a<red> d<red>} \\
            }
        }
    }
    & & \\ \\

    \theutterance \stepcounter{utterance}  
    & & & \multicolumn{4}{p{0.6\linewidth}}{
        \cellcolor[rgb]{0.9,0.9,0.9}{
            \makecell[{{p{\linewidth}}}]{
                \texttt{\tiny{[P1$\langle$GM]}}
                \texttt{gissning\_återkoppling: s<grön> å<gul> r<grön> a<röd> d<röd>} \\
\\ 
\texttt{Ge ditt svar endast i detta format.} \\
\texttt{förklaring: detaljer} \\
\texttt{gissning: ord} \\
            }
        }
    }
     \\ \\

    \theutterance \stepcounter{utterance}  
    & \multicolumn{4}{p{0.6\linewidth}}{
        \cellcolor[rgb]{0.561,0.737,0.561}{
            \makecell[{{p{\linewidth}}}]{
                \texttt{\tiny{[P1$\rangle$GM]}}
                \texttt{förklaring: s och r är fasta på plats 1 och 3. å är gul, alltså finns den i ordet men inte på plats 2. a och d är uteslutna. eftersom position 2 är den enda lediga platsen för å (platser 1,3,4,5 är upptagna av s,r eller uteslutna för å), måste å stå på position 4. jag testar nu "syrås" som fyller i å på plats 4 och testar y på plats 2 samt s på plats 5.} \\
\texttt{gissning: syrås} \\
            }
        }
    }
    & & \\ \\

    \theutterance \stepcounter{utterance}  
    & & & \multicolumn{2}{p{0.3\linewidth}}{
        \cellcolor[rgb]{0.9,0.9,0.9}{
            \makecell[{{p{\linewidth}}}]{
                \texttt{\tiny{[GM$|$GM]}}
                \texttt{Error: The guessed word is not a valid word for this game.} \\
            }
        }
    }
    & & \\ \\

    \theutterance \stepcounter{utterance}  
    & & & \multicolumn{4}{p{0.6\linewidth}}{
        \cellcolor[rgb]{0.9,0.9,0.9}{
            \makecell[{{p{\linewidth}}}]{
                \texttt{\tiny{[P1$\langle$GM]}}
                \texttt{Din gissning är inte ett giltigt ord för det här spelet. Vänligen försök igen.} \\
\\ 
\texttt{Ge ditt svar endast i detta format.} \\
\texttt{förklaring: detaljer} \\
\texttt{gissning: ord} \\
            }
        }
    }
     \\ \\

    \theutterance \stepcounter{utterance}  
    & \multicolumn{4}{p{0.6\linewidth}}{
        \cellcolor[rgb]{0.561,0.737,0.561}{
            \makecell[{{p{\linewidth}}}]{
                \texttt{\tiny{[P1$\rangle$GM]}}
                \texttt{förklaring: s och r är fasta på plats 1 och 3. å måste ligga på plats 4 (eftersom den är gul på plats 2 och platser 1,3,5 är otillgängliga). a, d, e, i, l, m, n, o, t, u är uteslutna. jag testar nu "skråv" som har k på plats 2, å på plats 4 och v på plats 5.} \\
\texttt{gissning: skråv} \\
            }
        }
    }
    & & \\ \\

    \theutterance \stepcounter{utterance}  
    & & & \multicolumn{2}{p{0.3\linewidth}}{
        \cellcolor[rgb]{0.9,0.9,0.9}{
            \makecell[{{p{\linewidth}}}]{
                \texttt{\tiny{[GM$|$GM]}}
                \texttt{Error: The guessed word is not a valid word for this game.} \\
            }
        }
    }
    & & \\ \\

    \theutterance \stepcounter{utterance}  
    & & & \multicolumn{4}{p{0.6\linewidth}}{
        \cellcolor[rgb]{0.9,0.9,0.9}{
            \makecell[{{p{\linewidth}}}]{
                \texttt{\tiny{[P1$\langle$GM]}}
                \texttt{Din gissning är inte ett giltigt ord för det här spelet. Vänligen försök igen.} \\
\\ 
\texttt{Ge ditt svar endast i detta format.} \\
\texttt{förklaring: detaljer} \\
\texttt{gissning: ord} \\
            }
        }
    }
     \\ \\

    \theutterance \stepcounter{utterance}  
    & \multicolumn{4}{p{0.6\linewidth}}{
        \cellcolor[rgb]{0.561,0.737,0.561}{
            \makecell[{{p{\linewidth}}}]{
                \texttt{\tiny{[P1$\rangle$GM]}}
                \texttt{förklaring: s och r är fasta på plats 1 och 3. å måste ligga på plats 4. a, d, e, i, l, m, n, o, t, u är uteslutna. jag testar nu "språng" som har p på plats 2, å på plats 4 och ng på plats 5 (giltigt ord).} \\
\texttt{gissning: språng} \\
            }
        }
    }
    & & \\ \\

    \theutterance \stepcounter{utterance}  
    & & & \multicolumn{2}{p{0.3\linewidth}}{
        \cellcolor[rgb]{0.9,0.9,0.9}{
            \makecell[{{p{\linewidth}}}]{
                \texttt{\tiny{[GM$|$GM]}}
                \texttt{Error: The length of the guessed word is not 5.} \\
            }
        }
    }
    & & \\ \\

    \theutterance \stepcounter{utterance}  
    & & & \multicolumn{2}{p{0.3\linewidth}}{
        \cellcolor[rgb]{0.9,0.9,0.9}{
            \makecell[{{p{\linewidth}}}]{
                \texttt{\tiny{[GM$|$GM]}}
                \texttt{game\_result = ABORT} \\
            }
        }
    }
    & & \\ \\

\end{supertabular}
}

\end{document}
